\documentclass[aps,prx,twocolumn,superscriptaddress,amsmath,amssymb]{revtex4-2}

\usepackage{graphicx}
\usepackage{dcolumn}
\usepackage{bm}
\usepackage{hyperref}
\usepackage{physics}

\begin{document}

\preprint{APS/123-QED}

\title{Exact Clifford-Jones Isomorphism: Formal Verification of Majorana Spinor Algebra in Polarization Optics}

\author{Goutev}
\email{author@institute.edu}
\affiliation{Independent Research Program in Information Geometry and Topological Photonics}

\date{\today}

\begin{abstract}
The search for non-Abelian Majorana zero modes is fundamentally a search for macroscopic manifestations of Clifford algebraic symmetries. While topological superconductors remain the primary physical candidate, their cryogenic constraints limit macroscopic exploration. In this theoretical companion paper, we demonstrate a strict, exact mathematical isomorphism between the $SU(2)$ chiral spinor algebra of Majorana fermions and the classical Jones calculus of polarization optics. Using the Lean 4 proof assistant, we formally verify that the Iwasawa (KAN) decomposition of $SL(2, \mathbb{C})$ optical transfer matrices identically satisfies the Clifford symmetries required for non-Abelian braiding. We mathematically formalize the Exceptional Point (EP) as a topological Drazin defect, proving that the geometric multiplicity collapse perfectly models the Majorana boundary condition. Furthermore, we reveal that linear polarizations constitute a zero-spin equator ($S_3 = 0$), acting as an exact optical analogue to the spin-singlet state, which strictly requires chirality to activate nonlinear self-induced gauge fields. This formally verified framework bridges deep geometric topology with classical optics, demonstrating that the topological symmetries of quantum matter are natively accessible in room-temperature non-Hermitian photonics.
\end{abstract}

\maketitle

\section{I. Introduction}
The realization of Majorana zero modes and their non-Abelian exchange statistics is typically treated as a complex many-body condensed matter problem. However, at its mathematical core, the physics of topological quantum computation is governed purely by geometry and the representation theory of Lie groups. Specifically, it relies on the symmetries of the Clifford algebra and the braiding of $SU(2)$ spinors.

Historically, classical optics and quantum topology have been treated as distinct mathematical regimes. The Jones calculus, developed in 1941 to describe classical polarization through birefringent crystals, relies on $2 \times 2$ complex matrices acting on two-component vectors. Remarkably, this formal structure is mathematically isomorphic to the $SU(2)$ spinor representation of the vacuum. 

In this work, we present a formally verified proof (via the Lean 4 proof assistant) of the exact isomorphism between Majorana spinor algebras and non-Hermitian Jones calculus. By distilling the physics into pure algebra without axiomatic debt, we demonstrate that polarization optics natively embodies the geometric symmetries required for non-Abelian braiding.

\section{II. Clifford Algebra and $SU(2)$ Spinors}
The algebra of Majorana fermions emerges from the Clifford algebra of real vector spaces. The creation and annihilation operators of standard fermions are decomposed into purely real (Hermitian) components, generating a representation of $SO(2N)$. 

When modeling the topological defect of a 1D Kitaev chain, the boundary modes are described by $SU(2)$ chiral spinors. The basis states are defined by their helicity:
\begin{equation}
|R\rangle = \frac{1}{\sqrt{2}} \begin{pmatrix} 1 \\ i \end{pmatrix}, \quad |L\rangle = \frac{1}{\sqrt{2}} \begin{pmatrix} 1 \\ -i \end{pmatrix}.
\end{equation}
These states possess well-defined spin angular momentum eigenvalues of $+\hbar$ and $-\hbar$ along the propagation axis. To simulate Majorana physics, an architecture must natively support $SU(2)$ rotations, hyperbolic boosts (for particle-hole mixing), and chirality inversion.

\section{III. The KAN-Jones Isomorphism}
We prove that the macroscopic transfer matrix $M(k) \in SL(2, \mathbb{C})$ of a non-Hermitian photonic lattice maps identically to these requisite operations through its Iwasawa (KAN) decomposition:
\begin{equation}
M = K A N.
\end{equation}

\paragraph{Compact Rotation ($K$):} The $K$ component is an $SO(2)$ subgroup element, representing standard Hermitian rotation. In Jones calculus, this corresponds to phase accumulation without energy exchange, generated by intrinsic birefringence. This natively implements the unitary rotations of the spinor basis.

\paragraph{Hyperbolic Boost ($A$):} The $A$ component is a non-compact, positive-definite diagonal matrix. In our isomorphism, this maps to balanced optical dichroism (gain and loss). Mathematically, it operates as a hyperbolic boost, identically mirroring the Bogoliubov transformations that mix particle and hole states in a superconductor.

\paragraph{Nilpotent Defect ($N$):} The $N$ component is an upper-triangular shear matrix. It is the geometric generator of the Exceptional Point (EP) and the heart of the topological trap.

\section{IV. The Defective Geometry of the Exceptional Point}
The most profound intersection of these fields occurs at the Exceptional Point. In standard Hermitian quantum mechanics, matrices are always diagonalizable. However, by balancing the Hermitian birefringence against the anti-Hermitian dichroism, the system's Hamiltonian coalesces into a defective matrix:
\begin{equation}
N(x) = \begin{pmatrix} 1 & x \\ 0 & 1 \end{pmatrix}.
\end{equation}
When $x \neq 0$, the characteristic polynomial is $(\lambda-1)^2$, yielding an algebraic multiplicity of 2. However, the eigenvector equation $(N-I)v = 0$ strictly requires the orthogonal component to vanish ($v_y = 0$). 

This represents a physical collapse of dimensionality—the geometric multiplicity strictly reduces to 1. In Lean 4, we have formally verified this collapse (\texttt{ep\_dimensional\_collapse}). This defective geometry acts as an algebraic "black hole" for polarization, flawlessly simulating the trapping mechanism of a topological defect at the boundary of a condensed matter system.

\section{V. The Zero-Spin Equator and Nonlinear Activation}
To complete the universal gate set, the topological defects must be braided. In optics, non-Abelian exchange requires photon interaction, synthesized via the circular Kerr effect ($\chi^{(3)}$).

The topological framework reveals a deep symmetry regarding linear polarizations. A linear state, such as horizontal polarization, is a perfect superposition of opposite chiralities:
\begin{equation}
|H\rangle = \frac{1}{\sqrt{2}}(|R\rangle + |L\rangle).
\end{equation}
On the Poincaré sphere, $|H\rangle$ resides exactly on the equator, possessing zero net spin angular momentum ($S_3 = 0$). It is the optical manifestation of a spin-singlet vacuum state.

Because the Kerr interaction Hamiltonian is proportional to the axial spin field, $H_{\mathrm{int}} \propto \chi^{(3)} S_3^{(1)} \otimes S_3^{(2)}$, states residing on the zero-spin equator experience no non-Abelian phase evolution. The nonlinear braiding is exclusively activated when the state is pushed off the equator into chirality ($S_3 \neq 0$). This zero-spin equator provides a naturally protected topological activation threshold, guaranteeing that Abelian noise traversing the linear basis cannot accidentally trigger a logical braid operation.

\section{VI. Formal Verification in Lean 4}
The rigor of this isomorphism is not reliant on physical approximations but on formal mathematical proofs. Using Lean 4, we have axiomatically mapped the $SU(2)$ chiral mechanics, the $SL(2, \mathbb{C})$ KAN decomposition, and the $\operatorname{Tr}(A_k) \equiv 0$ connection into a strictly verified codebase. 

The verification of chirality inversion via the half-wave plate ($HWP |R\rangle = -i|L\rangle$), the dimensional collapse of the EP, and the exact compilation of the braid matrix $B = e^{-i\pi\sigma_x/4}$ ensures zero axiomatic debt. The underlying mathematics of classical polarization is structurally identical to that of topological quantum superconductors.

\section{VII. Conclusion}
We have established a rigorously verified theoretical bridge between classical Jones calculus and the non-Abelian spinor algebra of topological quantum computing. By demonstrating that non-Hermitian optical phenomena natively implement the exact Clifford symmetries required for Majorana braiding, we provide a unified geometric framework. This isomorphism proves that the macroscopic symmetries of polarization optics are fundamentally identical to the microscopic topological structures of exotic quantum matter.

\end{document}
